\section{The exhibit}\label{sec:exhibit}

We test-drove the pipeline on one goal: explanation of multi-objective reasoning in SE, 2021--2026.
The numbers below are the pipeline's own outputs, kept in the replication package.

\textbf{Finding 1: the canon is paywalled.}
Of 54 papers selected for reading (43 recent, above the citation knee; 11 classics, found by snowballing), only 22 (41\%) had fetchable open-access PDFs: 47\% of the recents, 18\% of the classics.

\textbf{Finding 2: abstracts under-report.}
Coding the 22 full texts against their abstracts (Table~\ref{tab:flips}): abstracts earned 9 of the 21 SE flags, 3 of the 10 multi-objective flags, 8 of 18 explanation flags.
Roughly half the topic signal is invisible to a reader of abstracts alone.

\textbf{Finding 3: binary matching lies in the other direction.}
Any-match coding over full text flagged 21 of 22 papers with every flag at once; a ten-page PDF mentions everything somewhere.
The frequency threshold restores discrimination, and it discriminates most on the multi-objective flag (11 of 22 papers flip): half this corpus mentions optimization, few work on it.
Mentions are not methods.

\begin{table}
\caption{Per-flag coding of 22 papers: abstract versus
full text, binary versus thresholded (thr = at least
0.5 matches per 1000 words).}
\label{tab:flips}
\small
\begin{tabular}{lrrrr}
\toprule
flag & abs & full-bin & full-thr & flips bin$\to$thr\\
\midrule
software task & 9 & 22 & 21 & 1\\
multi-objective & 3 & 21 & 10 & 11\\
explanation & 8 & 22 & 18 & 4\\
benchmarking & 5 & 22 & 20 & 2\\
\bottomrule
\end{tabular}
\end{table}
